% !TeX root = ../main.tex

\chapter*{导言}
\addcontentsline{toc}{chapter}{导言}

整数维数与分数维数集合的几何测度论自本世纪初以来一直由纯数学家们发展.近来，分形集合在科学中的重要性急剧上升.Mandelbrot（\citeyear{mandelbrot1975,mandelbrot1977,mandelbrot1982}） 开创了用分形集合为从分子到天文尺度的各种科学现象建模的先河，例如：粒子的布朗运动、流体中的湍流、植物的生长、地理上的海岸线与地面、宇宙中星系的分布，甚至股票市场价格的波动.分数维数集合也出现在纯数学的诸多分支中，例如数论与非线性微分方程.更多的例子在 Mandelbrot 的科学性、哲学性与图示性的论述中有所描述.于是，起源于纯数学的概念在科学中找到了大量应用，而这些应用又为数学家提供了进一步问题的丰富源泉.本专著主要关心这类集合的几何理论，而非其应用.

\newterm{fractal}{分形}{fractal}（fractal）一词由 Mandelbrot（\citeyear{mandelbrot1975}）从意为``破碎的''拉丁词 \emph{fractus} 导出，他给出了分形的一个``尝试性定义'':Hausdorff 维数严格大于拓扑维数的集合；但他同时指出这一定义并不令人满意，因为它排除了某些高度不规则的集合，而这些集合按其精神显然应当视为分形.此前数学家们曾以多种方式称呼这类集合——``分数维数集合''、``Hausdorff 测度集合''、``具有精细结构的集合''或``不规则集合''.对这类集合的任何严格研究也必须包含对那些拓扑维数与 Hausdorff 维数相等的集合的考察，哪怕只是为了将它们从讨论中排除出去.因此，我在本书中收录这类正则集合（光滑曲线与曲面等）时并不感到抱歉.

为了推广点、曲线、曲面分别具有维数 $0$、$1$、$2$ 的观念，人们提出了许多估计``细的''或``高度不规则''集合的``大小''或``维数''的方法.由 Hausdorff 测度定义的 Hausdorff 维数，从数学家的观点看具有压倒性的优点：Hausdorff 测度是一个测度（即在两两不交集合的可数族上具有可加性）.

遗憾的是，即便是相对简单的集合，其 Hausdorff 测度与维数也可能很难计算；特别是这些量的下界往往不易得到.这已成为物理应用中的一个相当大的缺陷，并导致 Hausdorff 维数定义的若干变体被采用，其中某些变体的采用甚至是无意的.

\textcite{hurewiczwallman1941}、\textcite{kahane1976}、\textcite{mandelbrot1982}（第 39 节）与 Tricot（\citeyear{tricot1981,tricot1982}） 综述并比较了这些 Hausdorff 维数的替代定义.它们包括熵（见 \textcite{hawkes1974}）、相似维数（见 \textcite{mandelbrot1982}），以及 \textcite{johnsonrogers1982} 的局部维数与局部测度.基于任何这样的定义都不难写出一本类似性质的书，但 Hausdorff 测度与维数无疑是研究得最广泛、使用得最广泛的.

通过定义外测度把区间长度的概念推广到更复杂的实数集合，这一想法出现得出人意料地晚.\textcite{borel1895} 利用测度估计集合的大小，从而构造了某些病态函数.这些思想被 \textcite{lebesgue1904} 采纳，成为其积分理论的基础概念.\textcite{caratheodory1914} 引入了更一般的 ``Carathéodory 外测度''.特别地，他在 $n$ 维 Euclid 空间中定义了 ``$1$ 维''或``线性''测度，并指出 $s$ 维测度对其他整数 $s$ 也可以类似地定义.\textcite{hausdorff1919} 指出 Carathéodory 的定义对非整数 $s$ 同样有价值.他通过证明 Cantor 著名的``中三分之一''集在 $s=\log 2/\log 3=0.6309\ldots$ 时具有正而有限的 $s$ 维测度来说明这一点.分数维数集合的概念由此诞生，而 Hausdorff 之名也被用于称呼相应的维数与测度.

从那时起，关于 Hausdorff 测度及 Hausdorff 可测集合几何的大量结果被发现.\textcite{rogers1970} 的著作对内在的测度论给出了出色的论述，而 \textcite{federer1969} 的学术巨著则给出了测度几何的一个非常一般的处理，并延伸到曲面面积与同调积分理论的问题.

关于 Hausdorff 测度及其几何的许多工作归功于 Besicovitch，读者在本书中将反复遇到他的名字.事实上，多年来几乎所有关于 Hausdorff 测度的已发表工作都带有他的名字，其中大量工作包含了极为精巧的论证.近年来他的学生们又做出了许多重要的进一步贡献.\textcite{burkill1971} 与 \textcite{taylor1975} 的讣告文章可以让读者对 Besicovitch 在这一学科中的影响之巨大有所体会.

显然，Besicovitch 曾打算写一本题为 \emph{The Geometry of Sets of Points}（点集的几何）的几何测度论著作，它在许多方面很可能与本卷相似.Besicovitch 于 1970 年去世后，Roy Davies 教授在 Helen Alderson 博士（1972 年去世）的协助下，整理出了可能构成 Besicovitch``第 1 章''的一个版本.这一章注定没有任何后续，但它对本卷的前面部分产生了相当大的影响.

在第 1 章中我们定义 Hausdorff 测度并研究其基本性质，说明在某些直截了当的情形下如何计算集合的 Hausdorff 维数与测度.

我们特别感兴趣的是维数为 $s$ 的 $s$-集，即 $s$ 维 Hausdorff 测度非零但有限的集合.对仅受如此弱条件限制的一类集合而言，其几何必然只能是一般点邻域性质的研究.因此接下来的三章讨论局部性质：集合在一点处的密度，以及集合绕其每一点的方向分布，即切线的存在性问题.分数维数集合与整数维数集合分别处理.分数维数集合必然是分形，但整数维数集合中正则的``曲线状''或``曲面状''集合与不规则的``分形''集合之间存在显著对照.

第 5 章引入强有力的网测度技术.这使我们能够证明任何具有无穷 $s$ 维 Hausdorff 测度的集合都包含一个 $s$-集，从而按要求把 $s$-集理论推广到更一般的集合.网测度还被用于研究集合的 Descartes 积的 Hausdorff 测度.

下一章处理集合到低维子空间的投影.在这一部分工作中，位势论方法作为直接几何方法的替代被引入.

第 7 章讨论 Kakeya 问题，即求在其中可以转动单位长线段的最小测度集合.文中讨论了若干变体，并通过对偶性与前一章的投影定理以及调和分析联系起来.

最后一章收录了若干杂例，用以说明书中前面遇到的一些思想.

参考文献列于书末并按年份引用.更多的文献目录见 Federer（\citeyear{federer1947,federer1969}）与 \textcite{rogers1970} 与 \textcite{mandelbrot1982}.
